\section*{第 \ref{ch:b2:atmospheric-thermodynamics} 章——大气热力学}

\begin{solution}{exo:b2:atmospheric-thermodynamics:1}
(a) $e = e_s(\qty{278.15}{K}) = \qty{8.73}{hPa}$，$e_s(\qty{283.15}{K}) = \qty{12.32}{hPa}$。
(b) $8.73/12.32 = \qty{70.9}{\percent}$。(c) $r = 0.622 \times 8.73/(850 - 8.73) =
\qty{6.45}{g/kg}$，则 $q = r/(1 + r) = \qty{6.41}{g/kg}$。
\end{solution}

\begin{solution}{exo:b2:atmospheric-thermodynamics:2}
(a) $\mathrm{d}\ln e_s/\mathrm{d}T = 5419/T^2$：每开尔文的变化率，
在 \qty{273.15}{K} 时为 \qty{7.26}{\percent}，在 \qty{293.15}{K} 时为 \qty{6.31}{\percent}，在
\qty{303.15}{K} 时为 \qty{5.90}{\percent}。(b) $e_s(\qty{20}{\celsius}) = \qty{23.66}{hPa}$ 与 $e_s(\qty{25}{\celsius})
= \qty{32.26}{hPa}$；比值为 1.363，对比 $1.063^5 = 1.357$。速率是
$\ln e_s$ 的导数，因此在 \qty{5}{K} 内的倍数为 $\exp(5 \times 0.0620) = 1.363$，其中 0.0620 为
中点 \qty{22.5}{\celsius} 处的速率；以 \qty{6.3}{\percent} 复利连乘五次，接近但并不精确。(c) $23.66/23.39 = 1.012$：该公式偏高 \qty{1.2}{\percent}，因为
它采用三相点的潜热，而潜热随温度升高而减小；
故公式斜率偏陡，误差随离
\qty{0}{\celsius} 的距离增大而增大。
\qty{0}{\celsius}.
\end{solution}

\begin{solution}{exo:b2:atmospheric-thermodynamics:3}
(a) $125 \times (30 - 14) = \qty{2000}{m}$。(b) $30 - 9.8 \times 2.0 = \qty{10.4}{\celsius}$。
(c) $14 - 1.8 \times 2.0 = \qty{10.4}{\celsius}$：温度与露点在
饱和层处相遇，这是必然的，该规律即此计算式。(d) $125 \times (33 - 14) = \qty{2375}{m}$：
云底随午后地面空气增温而抬升。
\end{solution}

\begin{solution}{exo:b2:atmospheric-thermodynamics:4}
(a) 自地面至约 \qty{908}{hPa}，露点曲线与温度曲线重合：
即最底层半公里几乎饱和，或接近饱和。(b) 在约 908 与 \qty{889}{hPa} 之间，
随气压减小，温度曲线向右转，转向更高的温度：
温度随高度升高。(c) LCL 约为 \qty{964}{hPa}，高出地面几十米；LFC
约为 \qty{730}{hPa}；EL 约为 \qty{235}{hPa}。(d) 逆温层顶处的露点
比温度低 \qty{11}{K}，在 \qty{850}{hPa} 处则低 \qty{20}{K}：暖而极干的
空气位于较冷而湿润的空气之上。其相对湿度低，说明其暖度并非凝结所得；
下沉的空气以干绝热递减率因压缩而增温，同时保持其\omterm{def:b2:atmospheric-thermodynamics:moist-air}{混合比}，
这正会形成此类层结。
\end{solution}

\begin{solution}{exo:b2:atmospheric-thermodynamics:5}
(a) $T_v = 303.15 \times (1 + 0.020/0.622)/1.020 = \qty{306.76}{K}$；近似式给出
$303.15 \times (1 + 0.61 \times 0.020) = \qty{306.85}{K}$。(b) $\rho = p/(R_dT_v) = 10^5/(287.0
\times 306.76) = \qty{1.1358}{kg/m^3}$，而干空气为 $10^5/(287.0 \times 303.15) =
\qty{1.1494}{kg/m^3}$：轻 \qty{1.2}{\percent}。(c) 由 $T_v - T =
\qty{3.6}{K}$ 可得。(d) 上升气块只需比周围空气轻即可；在温暖海面上，
空气中的水汽所增加的浮力相当于数开尔文，
即一个可观的温度超额。
\end{solution}

\begin{solution}{exo:b2:atmospheric-thermodynamics:6}
(a) $e_s$ = 6.11, 23.66 与 \qty{43.54}{hPa}；$r_s = 0.622\,e_s/(1000 - e_s)$ = 3.82, 15.07 与
\qty{28.31}{g/kg}；$\Gamma_s$ = 6.48, 4.20 与 \qty{3.44}{K/km}。(b) 干气块降温
\qty{9.8}{K}，饱和气块降温 \qty{4.2}{K}：饱和气块比干气块暖 \qty{5.6}{K}。
(c) 高处空气更冷，$r_s$ 更小，可供凝结的水汽所剩无几，
且 $\Gamma_s$ 向 $\Gamma_d$ 回归。
\end{solution}

\begin{solution}{exo:b2:atmospheric-thermodynamics:7}
(a) $e = e_s(\qty{291.65}{K}) = \qty{21.52}{hPa}$，$r = 0.622 \times 21.52/(971.1 - 21.52) =
\qty{14.1}{g/kg}$；$\theta = 292.15 \times (1000/971.1)^{0.2857} = \qty{294.6}{K}$；$\theta_e =
294.6 \exp[2.501 \times 10^6 \times 0.0141/(1004.5 \times 291.5)] = 294.6 \times 1.128 =
\qty{332.3}{K}$。(b) $e_s(\qty{264.25}{K}) = \qty{3.13}{hPa}$，$r_s = 0.622 \times 3.13/(500 -
3.13) = \qty{3.92}{g/kg}$；$\theta = 264.25 \times 2^{0.2857} = \qty{322.1}{K}$；该气块
已饱和，故 $T_{\text{LCL}} = T$，且 $\theta_e = 322.1 \exp[2.501 \times 10^6 \times
0.00392/(1004.5 \times 264.25)] = \qty{334.3}{K}$。(c) $\theta$ 升高了 \qty{27.5}{K}，
其增温来自所凝结的 \qty{10.2}{g/kg} 水汽的潜热；$\theta_e$ 变化了
\qty{0.6}{\percent}，处于其近似公式的精度范围内，故其守恒。
\end{solution}

\begin{solution}{exo:b2:atmospheric-thermodynamics:8}
(a) $\text{CAPE} = 287.0 \times 4 \times \ln(700/300) = 287.0 \times 4 \times 0.847 =
\qty{973}{J/kg}$。(b) $\sqrt{2 \times 973} = \qty{44.1}{m/s}$。(c) 在 1 与
\qty{1500}{J/kg} 之间：为正，但不大。(d) 两条
曲线之间水平宽度不变的带；由\cref{prop:b2:atmospheric-thermodynamics:areas}可知其面积正比于 $\int
\Delta T\,\mathrm{d}\ln p = 4 \ln(700/300)$，与等温线的倾斜无关。
\end{solution}

\begin{solution}{exo:b2:atmospheric-thermodynamics:9}
(a) $D_c = \sqrt{3 \times 0.61 \times 80^3/2.1}\,\unit{nm} = \qty{668}{nm} =
\qty{0.668}{\micro m}$；$s_c = 2A/(3D_c) = 2 \times 2.1/(3 \times 668) = 0.00210 =
\qty{0.210}{\percent}$。(b) 由 $s_c = \sqrt{4A^3/(27\kappa D_d^3)}$ 解出 $D_d$：$D_d =
\bigl(4A^3/(27\kappa s_c^2)\bigr)^{1/3}$，在 \qty{0.3}{\percent} 时为 \qty{63}{nm}，在
\qty{0.1}{\percent} 时为 \qty{131}{nm}。(c) 所有介于 63 与 \qty{131}{nm} 之间的颗粒现在也会活化：
液滴数增多，且在液态水量相同时，液滴也更小。
\end{solution}

\begin{solution}{exo:b2:atmospheric-thermodynamics:10}
(a) $w = \sqrt{2 \times 130} = \qty{16.1}{m/s}$，对近地面静止的空气而言不可能。
(b) 逆温层顶的位温为 $293.65 \times (1000/889.1)^{0.2857} =
\qty{303.7}{K}$；具有此 $\theta$ 的近地面气块，其 $T = 303.7 \times (971.1/1000)^{0.2857}
= \qty{301.2}{K}$，即 \qty{28.0}{\celsius}。(c) $125 \times (28.0 - 18.5) = \qty{1190}{m}$，
高出地面，高于旧的逆温层顶。(d) 锋面沿其坡面抬升空气；
以及高空扰动前方空气的缓慢上升，后者使盖层本身冷却。
\end{solution}

\begin{solution}{exo:b2:atmospheric-thermodynamics:11}
(a) 当 $r_s \to 0$ 时，分子趋于 $g$，分母趋于 $c_p$。(b) 写出
$L^2r_s/(R_vT^2) = c_p \times \bigl(Lr_s/(R_dT)\bigr) \times LR_d/(R_vTc_p)$，其中在
\qty{303}{K} 时 $LR_d/(R_vTc_p) = 5.1$：当 $Lr_s/(R_dT) \gg 1$ 时，亦有 $L^2r_s/(R_vT^2) \gg c_p$，
则 $\Gamma_s \to g\,(Lr_s/(R_dT))/(L^2r_s/(R_vT^2)) = gR_vT/(LR_d) = 9.81 \times 461.5 \times
303/(2.501 \times 10^6 \times 287.0) = \qty{1.91}{K/km}$。(c) 当 $T_d = T$ 时，
$\Gamma_{\text{dew}} = gR_vT/(LR_d)$，二者相同。在此极限下，潜热极大，以至于
气块几乎不能凝结其水汽而不至于升温远高于其饱和点：
即它保持\omterm{def:b2:atmospheric-thermodynamics:moist-air}{混合比}不变，沿等\omterm{def:b2:atmospheric-thermodynamics:moist-air}{混合比}线运动，沿该线
温度以 $\Gamma_{\text{dew}}$ 的速率下降。(d) 在 \qty{30}{\celsius} 与 \qty{1000}{hPa} 下，
$Lr_s/(R_dT) = 0.81$，不远大于 1，且 $L^2r_s/(R_vT^2) = 4180$，为 $c_p$ 的四倍：
$\Gamma_s = \qty{3.4}{K/km}$，仍远离 \qty{1.9}{K/km}。该极限并未达到。
\end{solution}

\begin{solution}{exo:b2:atmospheric-thermodynamics:12}
(a) $2.868/2.599 = 1.104$ 与 $1.257/1.032 = 1.218$：\qty{10.4}{\percent} 与
\qty{21.8}{\percent}。(b) 液滴在\omterm{def:b2:atmospheric-thermodynamics:microphysics}{过饱和}度仅零点几个百分点时即可活化，
所需\omterm{def:b2:atmospheric-thermodynamics:microphysics}{过饱和}度少几十倍。(c) 冰晶通过凝华迅速生长，它们从空气中夺取水汽，
使液滴上方空气低于饱和；液滴蒸发，其水汽转移至
冰晶（贝吉龙过程），直至云全部冰化。(d) 在 \qty{0}{\celsius} 时两者饱和\omterm{def:b2:atmospheric-thermodynamics:moist-air}{水汽压}相等：
饱和\omterm{def:b2:atmospheric-thermodynamics:moist-air}{水汽压}相等：两者交于三相点，
其比值在该点附近趋于一。
\end{solution}

\begin{solution}{pb:b2:atmospheric-thermodynamics:1}
\textbf{1.} $e = e_s(\qty{291.65}{K}) = \qty{21.5}{hPa}$.
\textbf{2.} $r = 0.622 \times 21.5/(971.1 - 21.5) = \qty{14.1}{g/kg}$.
\textbf{3.} $q = 14.1/1.0141 = \qty{13.9}{g/kg}$.
\textbf{4.} $e_s(\qty{292.15}{K}) = \qty{22.2}{hPa}$; $21.5/22.2 = \qty{97}{\percent}$.
\textbf{5.} $0.61 \times 0.0141 \times 292.15 = \qty{2.5}{K}$.
\textbf{6.} $125 \times 0.5 = \qty{63}{m}$; $p = 971.1 \exp[-9.81 \times 62.5/(287.0 \times 294.7)]
= \qty{964}{hPa}$.
\textbf{7.} $\theta = 292.15 \times (1000/971.1)^{0.2857} = \qty{294.6}{K}$.
\textbf{8.} $16.9 - 17.3 = \qty{-0.4}{K}$：更冷。
\textbf{9.} $(20.5 - 16.4)/(1102 - 924) = 4.1/178$：气温升高 \qty{23}{K/km}，
即递减率为 \qty{-23}{K/km}。逆温：绝对稳定，相当于一个盖子。
\textbf{10.} $15.4 - 20.5 = \qty{-5.1}{K}$.
\textbf{11.} $e_s(\qty{-1.8}{\celsius}) = \qty{5.4}{hPa}$ 与 $e_s(\qty{18.6}{\celsius}) =
\qty{21.7}{hPa}$：\qty{25}{\percent}。温暖干燥，可能是下沉的空气，经压缩而增温；
也可能是上风方向干燥地面加热的空气；它并非下方的湿润空气。
\textbf{12.} $18.6 + 18.4 = \qty{37.0}{K}$ 除以 $5725 - 1489 = \qty{4236}{m}$，\qty{8.73}{K/km}：介于 $\Gamma_s$ 与
$\Gamma_d$ 之间，\omterm{def:b2:atmospheric-thermodynamics:instability}{条件不稳定}。
\textbf{13.} $-18.4 - (-8.9) = \qty{-9.5}{K}$：极度不稳定（$-8$ 至 $-10$）。
\textbf{14.} 由地面向上：0、$-0.4$、$-0.2$、$-5.1$、$-4.9$、$+1.6$、$+9.5$、$+10.0$、
$+5.6$、$+1.5$、$-9.5$~\unit{K}。
\textbf{15.} 在 850 ($-4.9$) 与 \qty{700}{hPa} ($+1.6$) 之间，取分数 $4.9/6.5 = 0.754$：
$\ln p = \ln 850 + 0.754 \ln(700/850)$, $p = \qty{734}{hPa}$; $z = 1489 + 0.754 \times 1621 =
\qty{2711}{m}$，高出地面 \qty{2.37}{km}。
\textbf{16.} 在 250 ($+1.5$) 与 \qty{200}{hPa} ($-9.5$) 之间，取分数 $1.5/11.0 =
0.136$：$p = \qty{242.5}{hPa}$，$z = 10560 + 0.136 \times 1446 = \qty{10757}{m}$，约
\qty{200}{m}，低于报告的对流层顶（\qty{10954}{m}）。
\textbf{17.} 各段（平均 $\Delta T$ 乘以气压比的 $\ln$）：$0.80 \times 0.0474
= 0.038$; $5.55 \times 0.3365 = 1.867$; $9.75 \times 0.2231 = 2.176$; $7.80 \times 0.2877 =
2.244$；$3.55 \times 0.1823 = 0.647$；$0.75 \times 0.0305 = 0.023$；求和得 \qty{6.995}{K}。
$\text{CAPE} = 287.0 \times 6.995 = \qty{2010}{J/kg}$，取三位有效数字。
\textbf{18.} 介于 1500 与 \qty{2500}{J/kg} 之间：较大。
\textbf{19.} $-0.2 \times 0.0486 - 0.3 \times 0.0189 - 2.65 \times 0.0207 - 5.0 \times 0.0450 -
2.45 \times 0.1467 = -0.655$~\unit{K}; $\text{CIN} = 287.0 \times (-0.655) = \qty{-188}{J/kg}$.
\textbf{20.} 表中仅保留十一个层次：各层之间，梯形法则以
$\ln p$ 中的直线段代替曲线，漏掉了其弯曲。\qty{1}{hPa} 的网格则依循曲线。\omterm{def:b2:atmospheric-thermodynamics:virtual}{虚温}
计入了气块水汽的浮力，饱和气块中的水汽比其周围的干空气更多，
因此 CAPE 更大，CIN 更小。
\textbf{21.} $\sqrt{2 \times 2010} = \qty{63}{m/s}$.
\textbf{22.} $10757 - 2711 = \qty{8046}{m}$，按 $63.4/2 = \qty{31.7}{m/s}$ 计：\qty{254}{s}，约四
分钟。
\textbf{23.} 上升气流所携带凝结水的重量；卷入周围较干的空气
（夹卷），随液滴蒸发而使其冷却。
\textbf{24.} 气块理论给出由此空气供给的风暴上升气流速度约为
\textbf{\qty[text-series-to-math=true]{63}{m.s^{-1}}}，一个上界。
\end{solution}
